\documentclass[answers]{exam}

\usepackage{../../../mypackages}
\usepackage{../../../macros}

% Pour le tableau
\usepackage{array}
\usepackage{longtable}
\usepackage{multirow}
\usepackage[table]{xcolor}

% -----------------------------
% Variables (valeurs par défaut — écrasées par gen_projet.py)
% -----------------------------
\newcommand{\varPrenom}{Anaëlle}
\newcommand{\varNom}{Duvault}
\newcommand{\varClasse}{6ème}
\newcommand{\varAnticipationNote}{5}
\newcommand{\varAnticipationAppreciation}{Excellent sur l’anticipation. Tu as très bien compris l’importance de commencer tôt et surtout de laisser des traces de ton travail. On voit que tu t’es mise sur le devoir en avance et que tu avances régulièrement. C’est exactement l’état d’esprit attendu dans ce type de projet.}
\newcommand{\varInteractionNote}{4}
\newcommand{\varInteractionAppreciation}{Très bonnes interactions par email. Tes messages sont clairs, structurés, avec du contexte, des explications sur ce que tu as essayé et souvent des captures d’écran pour montrer où tu bloques. Cela rend les échanges très efficaces. En classe, tu pourrais parfois venir un peu plus me voir quand une question te bloque.}
\newcommand{\varAttitudeNote}{4}
\newcommand{\varAttitudeAppreciation}{Très bonne attitude et beaucoup de sérieux dans le travail. Tu es volontaire et impliquée. Ton axe de progression est plutôt dans la dynamique de classe : essaie d’aller davantage vers d’autres élèves, proposer ton aide quand tu as compris et travailler aussi avec des personnes avec qui tu as moins l’habitude d’être.}
\newcommand{\varFormulasNote}{4.5}
\newcommand{\varFormulasAppreciation}{Très bon travail sur les formules et la structure des feuilles. Les VLOOKUP sont bien utilisés, les liens entre cellules sont cohérents et les sources sont correctement renseignées. On voit une vraie progression dans ta compréhension de Google Sheets. Un axe d’amélioration important cependant serait d’aller plus loin dans l’interprétation des résultats : parfois, tu donnes le sentiment de ne pas saisir exactement ce que l'on fait, à quoi correspond le sens de chaque colonne dans laquelle on calcule des résultats, et à quoi correspondent les étapes intermédiaires. Ce sont des données / metriques pas évidentes à comprendre certes, mais il faut à chaque étape quand même se demander ce qu'on est en train de calculer, et pourquoi on le fait. Ca permet d'être encore plus lucide sur les formules qu'on est en train d'écrire.}
\newcommand{\varNoteFinale}{17.5}
\newcommand{\varAppreciationGenerale}{Excellent trimestre. Tu as très bien compris les attentes du projet : anticiper, tenir au courant de ton avancée et structurer ton travail. Les formules sont bien maîtrisées et ton attitude est très positive. Continue sur cette dynamique, c’est du très bon travail.}

\title{Projets Scientifiques -- Compte rendu}
\author{Évaluation de \varPrenom{} \varNom{}}
\date{Mars 2026}

% Mise en forme des solutions en bleu
\SolutionEmphasis{\color{blue}}
\renewcommand{\solutiontitle}{\noindent}

\definecolor{lavande}{RGB}{180,190,255}
\definecolor{noteRed}{RGB}{220,50,50}
\definecolor{notePaleRed}{RGB}{255,180,180}
\definecolor{noteYellow}{RGB}{255,230,100}
\definecolor{notePaleGreen}{RGB}{180,230,180}
\definecolor{noteDarkGreen}{RGB}{40,140,40}

% Colonne p{} alignée à gauche et qui wrappe
\newcolumntype{L}[1]{>{\raggedright\arraybackslash}p{#1}}

% -----------------------------
% Commande de coloration conditionnelle des notes
% -----------------------------
\newcommand{\noteCell}[1]{%
  \edef\argval{#1}%
  \def\defaultval{--}%
  \ifx\argval\defaultval
    #1%
  \else
    \pgfmathtruncatemacro{\noteLevel}{ceil(\argval)}%
    \ifcase\noteLevel
      % 0
      \cellcolor{noteRed}\textcolor{white}{\textbf{#1}}%
    \or
      % 0.5–1
      \cellcolor{noteRed}\textcolor{white}{\textbf{#1}}%
    \or
      % 1.5–2
      \cellcolor{notePaleRed}\textbf{#1}%
    \or
      % 2.5–3
      \cellcolor{noteYellow}\textbf{#1}%
    \or
      % 3.5–4
      \cellcolor{notePaleGreen}\textbf{#1}%
    \or
      % 4.5–5
      \cellcolor{noteDarkGreen}\textcolor{white}{\textbf{#1}}%
    \fi
  \fi
}

\begin{document}

\textbf{Collège Lycée Suger}
\hfill
\textbf{Projets Scientifiques} \\

\textbf{Année 2025-2026}
\hfill
\textbf{Classe de \varClasse{}} \par

{\let\newpage\relax\maketitle}

\begin{tcolorbox}[
  colframe=red!90!black,
  fonttitle=\bfseries,
  boxsep=4pt,
  sharp corners,
]
\begin{center}
\textbf{\textcolor{red}{Ce document doit être imprimé et collé dans le cahier de projets scientifiques.}}
\end{center}
\end{tcolorbox}

\subsection*{Note}

\renewcommand{\arraystretch}{1.7}

\noindent
\begin{longtable}{|L{0.14\textwidth}|L{0.57\textwidth}|>{\centering\arraybackslash}p{0.08\textwidth}|>{\centering\arraybackslash}p{0.1\textwidth}|}
\hline
\rowcolor{lavande}
\textbf{Catégorie} & \textbf{Appréciation} & \textbf{Note} & \textbf{Barème} \\[2pt]
\hline

Anticipation
  & \varAnticipationAppreciation{}
  & \noteCell{\varAnticipationNote}
  & 5 \\
\hline

Interaction avec l'enseignant
  & \varInteractionAppreciation{}
  & \noteCell{\varInteractionNote}
  & 5 \\
\hline

Attitude et effort
  & \varAttitudeAppreciation{}
  & \noteCell{\varAttitudeNote}
  & 5 \\
\hline

Formules et sources
  & \varFormulasAppreciation{}
  & \noteCell{\varFormulasNote}
  & 5 \\
\hline

% ---------------- Total ----------------
\rowcolor{green!35}
\textbf{Total} & & \textbf{\varNoteFinale{}} & \textbf{20} \\

\hline

\end{longtable}

\subsection*{Appréciation générale}

\begin{tcolorbox}[
  colback=blue!5!white,
  colframe=blue!70!black,
  title=Commentaire,
  fonttitle=\bfseries,
  boxsep=4pt,
  arc=4pt,
]
\varAppreciationGenerale{}
\end{tcolorbox}


\subsection*{Rappel -- Barème de notation}

\renewcommand{\arraystretch}{1.7}

\noindent
\begin{longtable}{|L{0.14\textwidth}|L{0.73\textwidth}|>{\centering\arraybackslash}p{0.04\textwidth}|}
\hline
\rowcolor{lavande}
\textbf{Catégorie} & \textbf{Description} & \textbf{/5} \\[2pt]
\hline

Anticipation
  & L'élève a-t-il travaillé en avance et planifié son travail ? On attend une progression régulière visible dans l'historique Google Sheets, pas un travail fait au dernier moment. Un travail clairement anticipé avec planification visible donne 5/5.
  & 5 \\
\hline

Interaction avec l'enseignant
  & Qualité des échanges par email et en classe. On attend des messages clairs, contextualisés, avec des captures d'écran si nécessaire et une explication de ce qui a été essayé. Les retours de l'enseignant doivent être pris en compte.
  & 5 \\
\hline

Attitude et effort
  & Implication générale pendant le projet : autonomie, persévérance, aide apportée aux camarades, capacité à solliciter de l'aide au bon moment. Un élève qui demande de l'aide intelligemment n'est pas pénalisé.
  & 5 \\
\hline

Formules et sources
  & Maîtrise des formules Google Sheets (VLOOKUP, SOMME, multiplication, AVERAGE\ldots), compréhension de la logique derrière les calculs et qualité des sources utilisées. Éviter les valeurs en dur et les IFERROR non justifiés.
  & 5 \\
\hline

% ---------------- Total ----------------
\rowcolor{green!35}
\textbf{Total} & \textbf{Note maximale possible} & \textbf{20} \\
\hline

\end{longtable}


\end{document}
